\documentclass[journal]{IEEEtran}

% Useful LaTeX packages
\usepackage{cite}
\usepackage{graphicx}
\usepackage{amsmath,amssymb,amsfonts}
\usepackage{textcomp}
\usepackage{xcolor}
\usepackage{url}
\usepackage[hidelinks]{hyperref}

% ── Core packages (inherited from progress_reportv2.tex) ──
\usepackage[utf8]{inputenc}
\usepackage{titlesec}
\usepackage[T1]{fontenc}
\usepackage[margin=0.7in]{geometry}

% ── Math ───────────────────────────────────────────────────────────────────
\usepackage{amsmath, amssymb, amsthm, amsfonts, mathtools, bm}

% ── Graphics, color, hyperlinks ────────────────────────────────────────────
\usepackage{graphicx, subcaption}
\usepackage{xcolor}
\usepackage{cleveref}

% ── Tables, lists, floats ──────────────────────────────────────────────────
\usepackage{booktabs, multirow, array}
\usepackage{caption}
\usepackage{enumitem}
\usepackage[most]{tcolorbox}

% ── Algorithms ─────────────────────────────────────────────────────────────
\usepackage{algorithm}
\usepackage{amsmath,amssymb,amsthm,mathtools,bm}
\usepackage{graphicx,xcolor,hyperref}
\usepackage{algorithm,algpseudocode}
\usepackage{booktabs,enumitem,caption,subcaption}
\usepackage{fancyhdr} % Added for headers
\usepackage{tikz}
\usetikzlibrary{positioning, arrows.meta, fit, backgrounds, calc, decorations.pathreplacing}
\usepackage{amsthm}
\usepackage[most]{tcolorbox}
\usepackage{xcolor}

\newtcolorbox{summarybox}{colframe=red!70, boxrule=1.5pt, arc=4pt, left=2pt, right=2pt, top=3pt, bottom=3pt, halign=center}

\hypersetup{colorlinks=true,linkcolor=blue!70!black,citecolor=green!50!black,urlcolor=blue!60!black}
\setlist{nosep}

% ── Macros ─────────────────────────────────────────────────────────────────
\newcommand{\R}{\mathbb{R}}
\newcommand{\E}{\mathbb{E}}
\newcommand{\Var}{\operatorname{Var}}
\newcommand{\KL}{\operatorname{KL}}
\newcommand{\Normal}{\mathcal{N}}
\newcommand{\eps}{\varepsilon}
\newcommand{\rdiff}{\rho_{\text{diff}}}
\newcommand{\rhist}{\rho_{\text{hist}}}


\setlength{\headsep}{10pt}

\theoremstyle{definition}   % roman body, bold header
\newtheorem{definition}{Definition}

\begin{document}
% --- Improved Title Block ---
\twocolumn[
  \begin{center}
    \vspace{1em}
    {\large \bfseries Diffusion Models for Limit Order Book Generation \par}
    \vspace{0.6em}
    {\large \itshape On the Tension Between Distributional Realism, Temporal Transferability, and Multivariate Coverage \par}
    \vspace{0.6em}
    {\normalsize Arturo Favara, Maxime Vallot, Nick Bernardini \par}
    \vspace{0.3em}
    {\small Massachusetts Institute of Technology (15.458), Spring 2026 \par}
    \vspace{1.5em}
  \end{center}
]

\pagestyle{fancy}
\fancyhf{} % Clear all headers and footers
\fancyhead[L]{\small A. Favara, M. Vallot, N. Bernardini} % Left: Authors
\fancyhead[C]{\small\thepage} % Center: Page
\fancyhead[R]{\small 15.458, Spring 2026} % Right: Course

% make the title area
\titlespacing*{\section}{0pt}{1.2ex}{0.6ex}
\titlespacing*{\subsection}{0pt}{1.0ex}{0.5ex}
\titlespacing*{\subsubsection}{0pt}{0.8ex}{0.4ex}
\titleformat{\paragraph}[runin]{\normalfont\bfseries}{}{0pt}{}
\titlespacing{\paragraph}{0pt}{0.5ex}{1em}

\thispagestyle{fancy} % Apply headers to title page
\setlength{\parskip}{0pt}
\setlist[itemize]{leftmargin=1.5em}
\setlist[enumerate]{leftmargin=1.5em}

% ══════════════════════════════════════════════════════════════════════════
\begin{abstract}
    
Lorem ipsum
\end{abstract}
% ~250 words, non-technical. Audience: senior manager who knows finance but
% not ML. Structure:
%   1. What we did (one sentence).
%   2. Why it matters (one sentence).
%   3. How we did it (two sentences: 20 model variants, 3 evaluation
%      frameworks, 3 stocks).
%   4. What we found (two sentences: three-axis tension; no architecture
%      wins on all three).
%   5. Practical implication (one sentence: practitioners must choose).
% No jargon. No acronyms. No math.

% ══════════════════════════════════════════════════════════════════════════
\section{Introduction and Motivation}
\label{sec:intro}
% ~2 pages. Narrative arc: PROBLEM → why GENERATIVE models → narrow to
% DIFFUSION → state CONTRIBUTION.

\subsection{The Backtesting Problem}
\label{sec:backtest_problem}
% Historical replay is non-responsive; the market doesn't react to the
% agent's actions. Quote Berti et al. 2025 (TRADES). Synthetic generation
% as the alternative.

\subsection{Why Diffusion Models for LOBs}
\label{sec:why_diffusion}
% Brief history: GANs (Li 2020, Coletta 2021-2022) → mode collapse.
% Diffusion (Ho 2020) → stable, mode-covering, flexible conditioning.
% TRADES (Berti 2025) → first diffusion LOB generator.
% Our starting point: fine-tune TRADES backbone on our pipeline.

\subsection{Contributions}
\label{sec:contributions}
% Enumerate clearly:
%   (i)   Systematic 20-architecture ablation across 5 design dimensions
%   (ii)  Three-way tension: distributional realism vs temporal
%         transferability vs multivariate coverage
%   (iii) Tension is structural to MSE-based diffusion on heavy-tailed data
%   (iv)  Agent-ranking metrics decoupled from all three quality dimensions
%   (v)   Cross-asset validation on INTC (large tick), TSLA (small tick),
%         SPY (ETF) testing INTC-specificity hypothesis.

% Optional: visual abstract teaser of Figure 10.1 (rank-bump chart).

\subsection{Paper Organization}
\label{sec:organization}
% One-sentence roadmap for each main section (total: 8), plus appendices.
% Example structure:
%   §2 describes the LOB data and its statistical properties.
%   §3 presents the diffusion architecture and the five design axes we vary.
%   §4 defines our three-axis evaluation framework.
%   §5 presents the main results, including a copula case study (§5.7) that illustrates the tension most vividly.
%   §6 examines agent-based stress testing.
%   §7 validates cross-asset.
%   §8 discusses implications.
%   Appendices provide detailed results, model specifications, and ablation analyses.

% ══════════════════════════════════════════════════════════════════════════
\section{Data}
\label{sec:data}
% ~3 pages, deep and pedagogical. A reader who has never seen TAQ should
% understand the LOB after this section.

\subsection{Limit Order Book Mechanics}
\label{sec:lob_mechanics}
% What is a LOB: orders, levels, bid/ask, spread, midprice.
% Event types: quote updates vs trades.
% INTC as a large-tick stock; spread pinned at 1 tick; 89% of events
% don't change the midprice. Why this is interesting AND pathological.
%
% FIGURE 3.1: "Anatomy of a Limit Order Book" (TikZ or matplotlib)

\subsection{Data Source and Feature Engineering}
\label{sec:data_source}
% NYSE TAQ via WRDS: quotes + trades, nanosecond resolution.
% LOBSTER reconstruction pipeline.
% The 8-feature tensor: delta_t, event_type (discrete), bid_dist, ask_dist,
% bid_sz, ask_sz, trade_sz, trade_sign, mid_return.
% Normalization (z-score from training set).
% Windowing (256-event non-overlapping) and stitching (10 windows →
% 2,560-event pseudo-days).
%
% Closing paragraph: brief look at the empirical feature-correlation
% structure as orientation for the reader (heavy diagonal,
% mechanical bid_dist↔ask_dist anti-correlation, faint OFI→return
% Spearman ≈ 0.07 — flagged here purely so it is recognisable when
% it returns load-bearing in §5.7 and §6).
%
% FIGURE 2.1: Spearman correlation heatmap of the 8 features
%             (was Fig 3.3 in earlier drafts).
%
% NOTE: The INTC zero-inflation pathology (point mass at mid_return = 0)
% used to be its own §2.3 here. We removed it because it doesn't earn
% its airtime until §5.7 (the copula story), which is the place the
% pathology actually does load-bearing work. The first time the reader
% sees a 0.89 frac-zero number is therefore in §5.7.1, where it
% immediately motivates the failure mode being analysed.

\subsection{Train / Validation / Holdout Split and Regime Labelling}
\label{sec:splits}
% 60 / 10 / 20 days. Regime labels: base, high_vol, toxic, thin
% (per VPIN + realized vol bucketing). Why regime conditioning matters:
% generator should produce stress scenarios, not average days.
%
% TABLE 3.1: Data split summary (days, events, period per ticker)

\subsection{Stylized Facts of Financial Returns}
\label{sec:stylized_facts}
% Cont (2001) checklist:
%   (i)   Heavy tails (kurtosis >> 3)
%   (ii)  Volatility clustering (slow ACF decay of |returns|)
%   (iii) Absence of return autocorrelation
%   (iv)  Trade-sign persistence (positive lag-1 ACF)
%   (v)   Volume-volatility correlation
%   (vi)  Leverage effect
% These become the G1 validation criteria (§4.4.1).
%
% FIGURE 3.4: Standalone real-INTC stylized fact panel
%   (return histogram, |return| ACF, return ACF, trade-sign ACF,
%    spread distribution, intraday volume U-shape).

% ══════════════════════════════════════════════════════════════════════════
\section{Methods --- Diffusion Model Architecture}
\label{sec:methods}
% ~4 pages. Start with DDPM basics (1 page, brief), then design space
% (the 5 axes), then evaluation framework. Every architectural choice
% gets an INTUITION before the math.

\subsection{Denoising Diffusion Probabilistic Models}
\label{sec:ddpm}
% Forward process: gradual Gaussian noise. Reverse: learned denoising.
% ε-prediction (Ho 2020). Training loss: MSE on predicted noise.
% Keep BRIEF — 1 page max with references for depth.
%
% FIGURE 4.1: Forward/reverse process diagram
% ALGORITHM 4.1: DDPM training pseudocode

\subsection{The Architectural Design Space}
\label{sec:design_space}
% Introduce the 5 axes as a TABLE first, then explain each:
%
% TABLE 4.1: Axes × Options × What it controls
%   - Prediction target:  ε, v, x_0
%   - Conditioning:       FiLM, AdaLN-Zero
%   - Preconditioning:    None, EDM
%   - Data transform:     z-score, copula, copula+dequantization
%   - Loss weighting:     Min-SNR γ ∈ {1, 5}, uniform

\subsubsection{Prediction Target: $\eps$ vs $v$}
\label{sec:pred_target}
% INTUITION first.
%   ε-pred: network learns the noise added → mode-covering, preserves
%           bulk correlations, can't commit to extreme values.
%   v-pred (Salimans & Ho 2022): predict √ᾱ ε − √(1-ᾱ) x_0 → uniform
%           loss variance, better tail calibration, breaks cross-feature
%           correlations.
%
% FIGURE 4.2: ε vs v target magnitude across noise levels.

\subsubsection{Conditioning: FiLM vs AdaLN-Zero}
\label{sec:conditioning}
% FiLM (Perez 2018): affine modulation γ, β from regime embedding.
% AdaLN-Zero (Peebles & Xie 2023, DiT): adaptive LayerNorm with zero
% init → 6 modulation slots per block × N blocks vs FiLM's 4 slots.
% Reference Work5's F.2 finding: FiLM modulation collapse.

\subsubsection{Preconditioning: EDM (Karras et al.\ 2022)}
\label{sec:edm}
% σ-dependent input/output scaling that normalizes loss magnitudes
% across noise levels. σ_data parameter calibration (v6/v7/v7_b sweep).
% INTUITION: keep network's operating range comfortable at every σ.

\subsubsection{Data Transform: z-score vs Copula}
\label{sec:copula_transform}
% z-score: subtract mean, divide by std. Standard, preserves relative
% structure.
% Copula (Sklar 1959): per-feature Gaussianization via empirical CDF
% + Φ^{-1}. Diffusion learns dependence; tails restored at sampling.
% Zero-inflation pathology: 89% point mass → CDF discontinuity →
% z-collapse. Dequantization fix (Theis 2016).
%
% FIGURE 4.3: Copula pipeline diagram, with/without dequantization.
% ALGORITHM 4.2: Copula transform with dequantization.

\subsubsection{Loss Weighting: Min-SNR-$\gamma$ (Hang et al.\ 2023)}
\label{sec:min_snr}
% Weight per-step loss by min(SNR(t), γ)/SNR(t).
% γ=5 (default) downweights low-noise; γ=1 over-corrects (falsified in
% Phase A, see Appendix).

\subsection{The 20 Trained Checkpoints}
\label{sec:checkpoints}
% TABLE 4.2: Full configuration matrix
%   columns: Model, Prediction, Conditioning, EDM, Transform, γ, Clip, Notes

\subsection{Evaluation Framework: The Three Axes}
\label{sec:eval_framework}

\subsubsection{Distributional Realism (G1)}
\label{sec:eval_g1}
% 16 stylized-fact checks (Cont 2001 + microstructure-specific).
% Pass/fail per check, aggregate count. Measures: "does synthetic data
% look like real market data?"

\subsubsection{Temporal Pattern Transferability (Predictive Score)}
\label{sec:eval_predictive}
% Following TRADES (Berti 2025): train LSTM on synth, test on real,
% measure MAE. Measures: "can downstream ML models learn useful patterns
% from synthetic data?"

\subsubsection{Multivariate Coverage (PCA Convex Hull)}
\label{sec:eval_pca}
% Per-tape feature vector → PCA → convex hull overlap with real.
% Measures: "does synthetic data cover the full range of market states?"

\subsubsection{Agent-Based Ranking ($\rdiff$) and Why It Decouples}
\label{sec:eval_agents}
% Avellaneda-Stoikov family agents; ρ_diff = Spearman correlation between
% synth and holdout agent rankings. Flag immediately: §7 will show this
% metric is decoupled from all three quality dimensions.
%
% FIGURE 4.4: Three-axis triangle diagram with example metrics per vertex.

% ══════════════════════════════════════════════════════════════════════════
\section{Results --- The Three-Axis Tension}
\label{sec:results}
% ~6-7 pages now (Copula Story merged in as §5.7). Three protagonists
% (v2, v5, v9) tell the main story; full 20-checkpoint ablation explains
% WHY; the copula case study (§5.7) is the most pedagogically rich
% instance of "fixing one axis breaks another".

% ── §5.0 PREAMBLE — define the three-axis table up-front ───────────────────
% Open §5 by stating, in 1-2 paragraphs:
%   (a) the three axes recapped from §3.4 (G1 = distributional realism,
%       Predictive = temporal pattern transferability, PCA = multivariate
%       state-space coverage),
%   (b) why we present them as a JOINT comparison rather than separate sweeps:
%       any single-axis "best generator" is misleading — the empirical claim
%       of this paper is that the axes are anti-correlated across the 20
%       trained checkpoints,
%   (c) anchor TABLE 5.1 (multi-axis comparison, 20 rows × the four metrics
%       we report most heavily: G1 / Pred×replay / PCA% / kurtosis / trade
%       fraction) and FIGURE 5.1 (hero ternary).
% A reader who only looks at §5's first page should already see the central
% tension. Everything else in §5 (subsections 5.1-5.7) is "why and how".
%
% TABLE 5.1 (HERE, top of §5): full multi-axis comparison
%             — produced by scripts/80_report_figures.py → tab_1_3_multi_axis.{csv,tex}
% FIGURE 5.1 (HERO): three-axis evaluation space
%             — produced by scripts/80_report_figures.py → fig_1_1_three_axis_scatter

\subsection{Three Representative Generators}
\label{sec:three_protagonists}
% v2 (ε, FiLM, z-score): "safe baseline" — good temporal patterns, weak
%                        tails.
% v5 (v, AdaLN, z-score): "heavy-tail champion" — kurtosis 1256, breaks
%                         correlations.
% v9 (ε, FiLM, copula+dequant): "distributional realist" — perfect trade
%                                fraction, best vol clustering, poor
%                                temporal transferability.
%
% TABLE 5.1: Three-protagonist headline metrics
% FIGURE 5.1 (HERO): Three-axis evaluation space (ternary plot)

\subsection{Distributional Realism Results}
\label{sec:res_g1}
% G1 16-check matrix. Return tail (log-log CDF). Trade-fraction bar.
% Per-feature violins for the three protagonists.
%
% FIGURE 5.2: G1 pass/fail heatmap (16 × 20)
% FIGURE 5.3: Return tail comparison (log-log CDF)
% FIGURE 5.4: Trade fraction bar chart (real reference line)
% TABLE 5.2: Full G1 results for v2, v5, v9 (with real)

\subsection{Temporal Transferability Results}
\label{sec:res_predictive}
% Predictive score across all 20 models.
% The surprise: scale-exploded models (v3_e9_noclip, std 4×) BEAT
% scale-matched models (v9, std 0.6×) on predictive MAE.
% ACF comparison showing temporal structure preservation.
%
% FIGURE 5.5: Predictive score bar chart
% FIGURE 5.6: ACF grid — real vs v2 vs v9
% TABLE 5.3: Predictive score for all 20 models

\subsection{Multivariate Coverage Results}
\label{sec:res_pca}
% PCA scatter + convex hulls. v8 (copula no dequant) → 85% coverage
% despite 8/15 G1 (coverage ≠ realism). v2_noclip → 0% coverage despite
% 4× predictive score.
%
% FIGURE 5.7: PCA coverage panels (4 representative models)

\subsection{The Anti-Correlation Across Architectures}
\label{sec:res_anticorrelation}
% Pairwise scatter of evaluation metrics across 20 checkpoints.
% Spearman ρ matrix between metrics. NOT cherry-picked; systematic
% across every architectural axis.
%
% FIGURE 5.8: Pairwise anti-correlation scatter (G1 vs Pred vs PCA)
% TABLE 5.4: Metric correlation matrix (Spearman)

\subsection{Architectural Ablation: What Drives What}
\label{sec:res_ablation}
% Per-axis effect on per-metric:
%   ε → v:           improves kurtosis, degrades predictive score
%   remove x0_clip:  dramatically improves predictive, destroys G1
%   add copula:      improves PCA + trade fraction, degrades predictive
%   FiLM → AdaLN:    negligible effect on most metrics
%
% FIGURE 5.9: Ablation impact heatmap (axes × metrics)
% FIGURE 5.10: The x0_clip effect (paired bar chart)

\subsection{Case Study: The Copula Story --- One Architecture, All Three Axes Move}
\label{sec:copula_story}
% Was its own §6 in earlier drafts; merged in here because it is the
% single most pedagogically rich instance of the §5 tension. Each of
% v8, v9, v2_remapped moves all three axes (G1, Pred, PCA) at once,
% and reading them in sequence makes the structural mechanism obvious.
% Approx 1-1.5 pages of the §5 budget.
%
% Also folds in the zero-inflation pathology that earlier drafts tried
% to introduce in §2: the empirical 89% point mass at mid_return = 0
% is the technical root cause of the copula failure, and earns its
% airtime here where it actually does load-bearing work.

\subsubsection{The Zero-Inflation Pathology (Why a Naive Copula Was Always Going to Fail)}
\label{sec:copula_zero_inflation}
% INTC mid_return: 89% exactly zero — point mass + continuous heavy tail.
% Conditional on nonzero: kurtosis ~56, std ~6.8e-5, skew ~-0.17.
% Why this matters HERE (not in §2): the empirical CDF used by the copula
% transform has a discontinuity at 0 that maps 89% of training data to a
% single z-value, so the diffusion model trained on that z-space gets
% degenerate inputs.
%
% FIGURE 5.11: Three-panel histogram (INTC mid_return, TSLA mid_return,
%              INTC mid_return | mid_return ≠ 0).
%              Was Fig 3.2 in earlier drafts.

\subsubsection{The Promise: Marginals from Inverse CDF, Joints from Diffusion}
\label{sec:copula_promise}
% Copula theory (Sklar): separate marginals from dependence. Predicted
% outcome: kurtosis 400-800, ρ_diff ≥ +0.40.

\subsubsection{The Failure (v8): z-Space Collapse}
\label{sec:copula_failure}
% mid_return z-mean = -2.55, z-std = 0.30. Root cause: CDF discontinuity
% at the 89% point mass (§5.7.1). Kurtosis: 2.6 (not 400-800).

\subsubsection{The Fix (v9): Dequantization}
\label{sec:copula_fix}
% Standard normalizing-flow technique (Theis 2016). Spread the point
% mass uniformly before CDF computation. Result: frac_zero recovered
% (0.99 vs real 0.89), kurtosis 129, trade fraction 4.9%.

\subsubsection{The Deeper Lesson: Spearman $\neq$ Economic Information}
\label{sec:copula_lesson}
% v9: excellent distributional realism, poor temporal transferability.
% Round-trip (Gaussianize → train → sample → inverse) adds noise that
% scrambles temporal micro-patterns the LSTM relies on.
% v2_remapped (post-hoc quantile mapping of v2): perfect Spearman
% preservation, 265× predictive score, ρ_diff flipped from +0.40 to
% −0.40. Spearman invariance is mathematically correct but
% economically insufficient.
%
% FIGURE 5.12: z-space before/after dequantization
% FIGURE 5.13: Zero-inflation CDF visualization (was Fig 6.2)
% FIGURE 5.14: v2 vs v9 metric comparison (grouped bars)

% ══════════════════════════════════════════════════════════════════════════
\section{Agent-Based Stress Testing}
\label{sec:agents}
% ~2 pages. Punchline: ρ_diff is decoupled from distributional quality
% for a specific, structural reason.

\subsection{Agent Design --- The Avellaneda-Stoikov Family}
\label{sec:agent_design}
% Brief AS framework (1 paragraph + reservation-price equation).
% A0 (constant), A1 (AS baseline), A2 (AS+OFI), A3 (AS+VPIN).
% Fill simulator: order matching against synthetic book.

\subsection{The Ranking Mechanism --- A1/A2 Are Indistinguishable}
\label{sec:ranking_mechanism}
% A1 and A2 economically indistinguishable (effect size < 0.3 on every
% metric). The "correct" A2 > A1 ordering driven by fill-rate
% differences, not return prediction. OFI→return R² = 0.0001 — the
% signal A2 uses is noise-level.

\subsection{$\rdiff$ Is Decoupled From Distributional Quality}
\label{sec:rdiff_decoupled}
% v8 (worst marginals): ρ_diff = +0.80
% v9 (best marginals):  ρ_diff = -0.80
% v2_remapped (perfect Spearman preservation): ρ_diff = -0.40
% Scatter shows no correlation with G1, predictive, or PCA.
%
% FIGURE 7.1: Agent PnL overlap (A0/A1/A2/A3 boxplots on truth)
% FIGURE 7.2: ρ_diff vs everything (4-panel scatter)
% TABLE 7.1: Agent decomposition — what actually separates A1 from A2

\subsection{Implications for Agent-Based Synthetic Stress Testing}
\label{sec:agent_implications}
% Agent-ranking evaluation requires LARGE, STABLE performance gaps.
% Our 4 agents are too similar. ρ_diff measures fill-timing noise,
% not generator quality. This is a finding, not a failure — it tells
% practitioners what NOT to do with synthetic stress tests.

% ══════════════════════════════════════════════════════════════════════════
\section{Cross-Asset Validation: INTC, TSLA, SPY}
\label{sec:cross_asset}
% ~1.5-2 pages. Tests whether the three-axis tension is structural or
% INTC-specific.

\subsection{Why Cross-Asset?}
\label{sec:why_cross}
% If the tension persists on TSLA (small tick) and SPY (ETF), it's
% structural. If not, it's an INTC zero-inflation artifact.

\subsection{TSLA --- Small-Tick Single Name}
\label{sec:tsla}
% Smaller relative tick, ~80% zero-return rate (vs INTC's 89%), more
% continuous returns. Tests whether the ε vs v kurtosis gap shrinks
% under less zero-inflation.
%
% TABLE 8.1: v2 vs v5 on TSLA — three-axis comparison

\subsection{SPY --- Index-Tracking ETF}
\label{sec:spy}
% Deep liquidity, tight spreads, index-arbitrage-driven price
% discovery. Tests whether the tension persists in a qualitatively
% different microstructure regime.
%
% TABLE 8.2: v2 vs v5 on SPY — three-axis comparison

\subsection{Cross-Asset Findings: Structural or Microstructure-Specific?}
\label{sec:cross_findings}
% Six data points (3 tickers × 2 architectures). Does v2 still beat v5
% on predictive across tickers? Does the trade-fraction divergence
% replicate? Plot the three-axis ternary with all 6 points.
%
% FIGURE 8.1: Cross-ticker three-axis ternary

% ══════════════════════════════════════════════════════════════════════════
\section{Discussion}
\label{sec:discussion}
% ~1.5 pages.

\subsection{The Three-Axis Tension as a Structural Result}
\label{sec:disc_structural}
% Why this happens: MSE is per-element, per-timestep. Cannot
% simultaneously optimize marginal shape AND cross-feature/temporal
% correlations. Connects to broader mode-covering vs mode-seeking
% diffusion literature.

\subsection{Practical Implications for Practitioners}
\label{sec:disc_practical}
% Choose generator by use case:
%   - Train downstream ML → prioritize predictive transferability
%     (use noclip / scale-exploded variants).
%   - Regulatory distributional stress test → prioritize G1
%     (use v2 or v5).
%   - Diverse scenario coverage → prioritize PCA (use copula).
% No one-size-fits-all answer.

\subsection{Limitations}
\label{sec:limitations}
% (i)   Single backbone (TRADES transformer); U-Net, SSM untested.
% (ii)  INTC-dominant (zero-inflation pathology); TSLA/SPY qualify.
% (iii) 20 holdout days insufficient for stable agent rankings.
% (iv)  Predictive score is LSTM-architecture-dependent.

\subsection{Future Work}
\label{sec:future_work}
% (i)   Autoregressive diffusion (event-by-event) might avoid
%       marginals-vs-joints tension.
% (ii)  Mixture-of-experts: ε-pred for bulk + v-pred head for tails.
% (iii) Direct distributional loss (MMD, Wasserstein) instead of MSE.
% (iv)  More agents with diverse strategies; longer holdout periods.

% ══════════════════════════════════════════════════════════════════════════
\section*{AI Disclosure}
\label{sec:ai_disclosure}
\addcontentsline{toc}{section}{AI Disclosure}
% "We used Claude (Anthropic) extensively as a research collaborator
% throughout this project for architectural design, experiment planning,
% diagnostic analysis, code generation, and report drafting. All final
% decisions, code execution, experimental results, and interpretations
% are our own. Claude's recommendations were frequently wrong (e.g., the
% copula kurtosis prediction) and we report negative results honestly."

% ══════════════════════════════════════════════════════════════════════════
\bibliographystyle{plain}
% \bibliography{references}
\section*{References}
\label{sec:references}
\addcontentsline{toc}{section}{References}
% Minimum citations:
%   Ho et al. 2020 (DDPM)
%   Song et al. 2020 (DDIM)
%   Salimans & Ho 2022 (v-prediction)
%   Karras et al. 2022 (EDM)
%   Peebles & Xie 2023 (DiT / AdaLN-Zero)
%   Perez et al. 2018 (FiLM)
%   Berti et al. 2025 (TRADES)
%   Cont 2001 (stylized facts)
%   Cont, Kukanov, Stoikov 2014 (OFI)
%   Avellaneda & Stoikov 2008 (market making)
%   Sklar 1959 (copula)
%   Theis et al. 2016 (dequantization for normalizing flows)
%   Hang et al. 2023 (Min-SNR weighting)
%   Nichol & Dhariwal 2021 (improved DDPM)
%   Li et al. 2020 (LOB GAN)
%   Coletta et al. 2021, 2022 (LOB GAN follow-ups)

% ══════════════════════════════════════════════════════════════════════════
\appendix
\section{Full G1 Value Matrix}
\label{app:g1_matrix}
% 16 × 20 actual-values table (not just pass/fail).

\section{All 20 Checkpoint Configurations}
\label{app:checkpoints}
% Detailed hyperparameters per model.

\section{Per-Feature Deep Diagnostic}
\label{app:deep_diagnostic}
% Output of scripts/60_deep_diagnostic.py for v2, v5, v9.

\section{Copula Transform Implementation}
\label{app:copula_impl}
% Dequantization algorithm; ε computation; per-feature treatment.

\section{Agent Implementation Details}
\label{app:agents_impl}
% Avellaneda-Stoikov parameters; fill simulator semantics; calibration.

\section{Predictive Score LSTM Details}
\label{app:lstm_details}
% Architecture, training curves, hyperparameter sensitivity.

\section{Additional Mid-Price Traces}
\label{app:traces}
% All 20 models, 2,560-event traces, normalized.

\section{Volume-Volatility Correlation Distributions}
\label{app:vol_vol}
% Per-tape KDE plots across model groups.

\section{Phase A Falsification --- Min-SNR $\gamma$ Misadventure}
\label{app:phaseA}
% Cautionary log: γ=1 hypothesis was directionally wrong; documented
% lesson on hyperparameter mathematical structure.

\section{Reproducer Commands}
\label{app:reproducer}
% Exact cluster commands (sbatch invocations) per checkpoint.

\end{document}